%======================================================================
% Section 2: Literature Landscape and Scholarly Context
%======================================================================
\section{Literature Landscape and Scholarly Context}
\label{sec:literature}

\subsection{The organising problem and the shape of the field}

The literature of generalized continuum wave mechanics did not accumulate gradually; it was triggered by a single anomaly and then organised itself around the successive attempts to remove it. The anomaly is the infinite propagation speed of Fourier heat conduction. For a body heated at its boundary, Fourier's law \eqref{eq:4.1} produces a temperature field that responds instantaneously everywhere, and the associated thermal wave number is unbounded. In 1948, in the proceedings of the Italian academy, Cattaneo \cite{Cattaneo1948} showed that the paradox disappears if heat flux is allowed a relaxation time; in 1958 Vernotte \cite{Vernotte1958} reached the same conclusion independently in the French academy proceedings. The two papers are, in structural terms, the same result arrived at from the same constraint by two separate routes, and they established the argument that the history of the field would follow: a physical inconsistency is identified, a constitutive modification is proposed that removes it, and the modified model is then shown to require further modification because something else has been lost.

That pattern --- identification, modification, cost --- is the organising theme of this review, and it is why the literature is best read by theory family rather than by paper. Each family below answers the same question in a different way, and each pays for the answer in a different currency: a second temperature, a second wave, a loss of ellipticity, a loss of uniqueness, a loss of local well-posedness. Sections~\ref{sec:thermoelastic} through~\ref{sec:advanced} take the families in turn; this section establishes how they arose, where the field's centre of gravity has moved, and what the geographic distribution of the work reveals about the structure of the problem.

\subsection{Four phases in the development of the field}

\paragraph{Phase I --- founding (1948--1972).} The initial period established the dichotomy that still organises the subject. The first branch modified only the thermal constitutive law and left the mechanical field untouched: this is the relaxation-time family, from Lord and Shulman \cite{Lord1967} onward, in which the heat equation becomes hyperbolic and a thermal wave appears. The second branch modified the set of state variables: Green and Lindsay \cite{Green1972} introduced two temperatures and two relaxation times, of which one family is dissipative and the other, under suitable symmetry conditions, conservative. Biot's classical coupled formulation \cite{Biot1956} is the common limit of both branches, and remains the benchmark against which every later model is judged.

\paragraph{Phase II --- consolidation and the experimental question (1972--1993).} During this period the question shifted from whether finite thermal propagation speed is admissible to whether it is observable. The second-sound literature supplied the decisive experimental context: Guyer and McCracken \cite{Guyer1990} established that the thermal modes of a solid can, at low temperatures, propagate as coherent waves rather than diffuse, which converts a modelling question into a physical one. The Green--Naghdi construction \cite{Green1991,Green1993} then introduced thermal displacement as a state variable, producing a type II theory in which the thermal field is conservative and a type III theory in which it is dissipative --- a distinction that matters because only the conservative variant admits a Hamiltonian structure and hence a well-defined energy. The synthesis of this phase is the review literature: Chandrasekharaiah's articles on hyperbolic thermoelasticity \cite{Chandrasekharaiah1986,Chandrasekharaiah1998}, Joseph and Preziosi's account of heat waves from kinetic theory to continuum mechanics \cite{Joseph1989,Joseph1990}, and the monographs by Nowacki \cite{Nowacki1975} and Achenbach \cite{Achenbach1973}.

\paragraph{Phase III --- the microscale turn (1995--2011).} Phase II assumed that the Fourier paradox was a matter of taking a limit too early. Phase III questioned the premise, and asked instead what a thermal wave should look like when the material has microstructure. Tzou's dual-phase-lag \cite{Tzou1995} and triple-phase-lag \cite{Tzou1997} models answer this by separating the lag associated with the flux from the lag associated with the temperature gradient, so that the two are not forced to coincide; the resulting hierarchy of thermal modes is the clearest demonstration in the literature that the thermal law and the mechanical law can be generalised at different rates. In parallel, Eringen's nonlocal elasticity \cite{Eringen1972,Eringen1974,Eringen1983} and the strain-gradient formulations \cite{Askes2006,Askes2011,Polizzotto2001} transferred the same logic to the mechanical side, and the size-dependent wave speeds measured and computed for nanostructures \cite{Maranganti2006,Maranganti2007} supplied the corresponding experimental evidence.

\paragraph{Phase IV --- memory, fractional operators and multiphysics (2011--present).} The most recent phase has replaced the single relaxation time with an entire memory kernel, and has replaced the single field with a coupled system. Memory-dependent derivatives \cite{Wang2011} generalise the exponential kernel into an arbitrary weighting function, which makes the memory structure a modelling choice rather than a parameter; fractional thermoelasticity \cite{Caputo1967,Gorenflo1997,Luchko2010,Povstenko2004,Povstenko2007,Povstenko2009,Povstenko2011,Povstenko2016} replaces the derivative order with a real exponent in $(0,2)$, and thereby interpolates the diffusive and wave-like regimes continuously rather than discretely; and the coupling of the thermal field to electrical, porous, micro-rotation and carrier fields \cite{Iesan1986,Iesan2013,Cowin1983,Nunziato1979,Chandrasekharaiah1996,Marotti2013,Sharma2022,Abouelregal2023} converts single-field wave problems into multiphysics ones in which the number of modes, and hence the size of the interface system, grows with the number of couplings.

\subsection{International structure of the field}

The development of generalized thermoelasticity was international from the outset, and the pattern of its internationalization follows the pattern of the science rather than that of any one school. The founding hyperbolic-conduction results appeared in the Italian and French academy proceedings \cite{Cattaneo1948,Vernotte1958}, and the consolidation of the subject in the 1970s occurred in the American journals of applied mechanics \cite{Biot1956,Lord1967,Green1972,Green1991,Green1993,Tzou1995}. The micropolar extension of the theory --- which adds micro-rotational degrees of freedom to a material with internal structure --- was developed across the Turkish, German and American continuum-mechanics communities and published in the engineering-science journals of that period \cite{Eringen1966a,Eringen1968b,Eringen1972}; the anisotropic extension followed from Canadian and Egyptian work on direction-dependent thermoelastic response \cite{Dhaliwal1980,Sharma2022}. The fractional-thermoelastic programme of the 2000s and 2010s was carried principally from the Ukrainian school, whose contributions span the time-fractional heat equation, the general fractional formulation and its convergence properties \cite{Povstenko2004,Povstenko2007,Povstenko2009,Povstenko2011,Povstenko2016}; the memory-dependent formulation was developed in China \cite{Wang2011}, and the mathematical apparatus of fractional viscoelasticity that underwrites it was placed on a rigorous footing in the Italian school \cite{Gorenflo1997,Luchko2010}. In the most recent period the activity has become genuinely multinational, with contributions from the Egyptian and Saudi groups on fractional and memory-based thermoelasticity and their interface problems \cite{Abouelregal2023,Sharma2022}, from Italian mathematical studies of the variational and phase-transition structure of multiphysics continua \cite{Marotti2013}, from Spanish work on well-posedness and uniqueness in nonlocal and fractional thermoelasticity \cite{Quintanilla2019}, and from an Indian group whose recent series addresses piezothermoelastic and memory-dependent wave problems, energy ratios at material interfaces and heat-map representations of coupled-field response \cite{Barak2023,Barak2023IJP,Barak2023JTS,Barak2023MMMS,Gupta2023Energy,Gupta2023IJNMF,Kumar2023PhysScr,Gupta2026HeatMap}.

This distribution is not incidental. Each of the centres identified above entered the field through a different scientific interest --- heat conduction in metals, micropolar fluid mechanics, anomalous diffusion in geophysics, piezoelectricity in smart materials --- and each imported the constitutive modification appropriate to that interest. The result is a literature in which the same physical question, the admissibility of a propagating thermal mode, is attacked with structurally different devices: a relaxation time, a second temperature, a thermal displacement, a phase lag, a memory kernel, a fractional order. A review organised around papers would obscure this, because the same device reappears in every centre with different notation. A review organised around the devices --- which is what this review attempts --- shows that the field is not a sequence of independent advances but a single continuing negotiation between three requirements that are individually necessary and jointly expensive: finite propagation speed, thermodynamic admissibility, and well-posedness.

\subsection{The literature corpus and the verification protocol}

The scholarship summarised in this review rests on a deliberately constructed corpus rather than on an opportunistic citation list. The corpus was assembled in three layers. The first layer is a set of foundational works identified by tracing the constitutive modifications backwards to their origin, and verified individually against the Crossref registry: the resolving metadata --- title, author list, journal, volume, issue, page range and year --- was taken from the registered record, not from a secondary citation list. The second layer is a set of review articles and monographs chosen to characterise the state of each theory family at the close of its formative period \cite{Chandrasekharaiah1986,Chandrasekharaiah1998,Shakeriaski2021,Nowacki1975,Achenbach1973}. The third layer is a set of primary research contributions selected to represent the current research front across coupled, nonlocal, fractional and memory-dependent families, including the recent interface-energy-ratio series \cite{Barak2023,Barak2023IJP,Barak2023JTS,Barak2023MMMS,Gupta2023Energy,Gupta2023IJNMF,Kumar2023PhysScr,Gupta2026HeatMap} and the fractional-thermoelastic work of the Egyptian group \cite{Abouelregal2023}.

The assembled database contains 95 records with bibliographic metadata, of which 48 carry a registered DOI that resolves to a record whose title and first author match the entry; the remaining 47 are pre-DOI works, monographs, or conference and report literature identified unambiguously by author, year, title and venue. Fifty-six of the works carry a verified DOI in the final bibliography, which is the set actually cited in the text of this review; the distinction between the size of the corpus and the size of the cited bibliography is deliberate, because the corpus is a discovery instrument and the bibliography is an assertion set. Every citation in this manuscript is to a work whose bibliographic record was checked against the registry or against the publisher's own record, and no work is cited on the strength of a secondary list alone.

Three consequences follow for how the remainder of the review should be read. First, coverage is theory-family-complete rather than application-complete: the corpus was assembled by constitutive structure, so an application is covered when a theory family that addresses it is covered, and not otherwise. Second, the evidence levels assigned in Section~\ref{sec:applications} distinguish between a mechanism that is experimentally established, a mechanism that is established but only modelled approximately, and a mechanism that remains speculative; the distinction is made in the review because the corpus contains all three kinds and because collapsing them would misrepresent the state of the field. Third, and most importantly for the treatment of individual contributions, the review is organised by scientific development and conceptual theme. Where two theories differ in their predictions, the review states what the difference is and what would distinguish them experimentally, rather than arranging the literature as a sequence of individual studies.

\subsection{Development axes and the gaps they expose}

The corpus can be classified along five axes --- physical domain, mathematical structure, modal content, coupling topology and validation regime --- and this classification is developed formally in Section~\ref{sec:taxonomy}. Its value at this stage is diagnostic: it shows where the field is dense and where it is thin. The dense regions are the single-field, one-dimensional, interface-focused problems of classical, relaxation-time and phase-lag thermoelasticity, for which a large body of modal and energy-ratio results exists. The thin regions are more revealing. There is very little work that closes the loop from constitutive novelty to experimental discrimination: the nonlocal and fractional families are theoretically mature, with well-established dispersion relations and modal structures, yet the number of studies that fit predictions to measured dispersion data remains small. The comparison of competing mechanisms for the same physical effect --- nonlocal versus gradient, exponential versus power-law memory, phase lag versus relaxation time --- is similarly under-represented, even though it is precisely in this comparison that the theory choices are testable. A third thin region is the treatment of interfaces in the presence of memory or nonlocality, where the constitutive history and the geometric matching condition must be reconciled, and where the energy partition is not determined by the local constitutive law.

These gaps are not defects of the individual contributions; they follow from the structure of the field, in which each advance has been made by extending a constitutive law and then solving the resulting modal problem, and the natural next step --- putting two such extensions in competition against an experiment --- has not yet become the standard mode of work. Section~\ref{sec:future} takes this diagnosis as the basis for its research directions.

\subsection{Existing reviews and the contribution of this article}

Reviews already exist for several individual theory families: nonlocal elasticity and its application to nanostructures \cite{Askes2006,Askes2011,Maranganti2007}, fractional thermoelasticity \cite{Chandrasekharaiah1998,Povstenko2016}, hyperbolic thermoelasticity and second sound \cite{Chandrasekharaiah1996,Joseph1989}, and generalized thermoelasticity with modified models \cite{Shakeriaski2021}. Their common structure is the one this review sets aside: an enumeration of theory families, each with its governing equations, its application domains and a list of its principal results, with the families treated in isolation and the transitions between them left implicit.

The present review differs in three respects. It is organised by the problems --- finite propagation speed, thermodynamic admissibility, well-posedness, scale dependence, interface partition of energy --- rather than by the theories, so that a reader interested in one problem can follow it across all the theories that address it. It states the cost of each modification alongside the benefit, which the standard format does not do, and which the literature itself rarely does: the second temperature, the loss of ellipticity, the loss of uniqueness of the initial-value problem, and the non-locality of the energy flux are each the price of something that was bought. And it separates the question of a mechanism's plausibility from the question of its verification, treating the evidence regime as an explicit axis of classification rather than as a matter of emphasis.

\begin{figure}[htbp]
\centering
\includegraphics[width=0.95\textwidth]{fig03-literature-landscape.pdf}
\caption{Theory family evolution in generalized continuum mechanics: the historical development from classical Fourier heat conduction, through the relaxation-time, two-temperature and thermal-displacement families, to the phase-lag, nonlocal, fractional, memory-dependent and coupled multiphysics formulations. The horizontal axis is the development of the thermal constitutive law; the vertical axis records the corresponding treatment of scale and internal structure (Section~\ref{sec:taxonomy} classifies the same corpus formally).}
\label{fig:literature_landscape}
\end{figure}

\subsection{Citation and verification principles}

Every equation presented in this review is selected for detailed exposition on four grounds. First, \textbf{verification status}: an equation is presented as a canonical result only if it is derived from first principles within the manuscript or taken from a primary source whose record has been checked against the registry, and the derivation is then reproduced with its assumptions stated. Second, \textbf{validity domain}: each equation is presented with its constitutive assumptions, material symmetry, geometry, boundary conditions, thermal theory, frequency convention, normalisation and approximation level made explicit, since the majority of apparent contradictions between published results are in fact differences of convention. Third, \textbf{open questions}: an equation that cannot be independently verified is identified as an open question rather than presented as settled, and is listed as such in the register of Appendix~A. Fourth, \textbf{minimality}: fewer precise equations are preferred over a larger number of redundant ones.

Where material errors are conclusively established in the literature, they are recorded transparently and without personal characterisation. Where uncertainty exists, it is acknowledged as an open question. No error allegations are made against any publication, and model-specific differences are never presented as mathematical contradictions without sufficient evidence.

% LocalWords: 
